\documentclass[11pt]{article}

\usepackage[utf8]{inputenc}
\usepackage[T1]{fontenc}
\usepackage{amsmath, amssymb, amsthm, mathtools}
\usepackage{microtype}
\usepackage[margin=1in]{geometry}
\usepackage{hyperref}
\usepackage{enumitem}
\usepackage{xcolor}
\usepackage{booktabs}

\hypersetup{
    colorlinks=true,
    linkcolor=blue!50!black,
    urlcolor=blue!70!black,
    citecolor=blue!50!black
}

\theoremstyle{plain}
\newtheorem{theorem}{Theorem}[section]
\newtheorem{proposition}[theorem]{Proposition}
\newtheorem{lemma}[theorem]{Lemma}

\theoremstyle{definition}
\newtheorem{definition}[theorem]{Definition}

\theoremstyle{remark}
\newtheorem{remark}[theorem]{Remark}

\title{Phase 4, Block 4 \\[0.3em]\Large Heston model: Andersen's Quadratic-Exponential (QE) scheme}
\author{Quant Finance Portfolio}
\date{\today}

\begin{document}
\maketitle
\tableofcontents
\bigskip

%============================================================
\section{Introduction and goals}
%============================================================

In Block 3 we implemented full-truncation Euler for Heston Monte Carlo. The scheme is robust, simple, and works correctly across all parameter regimes, but it carries a discretisation bias of order $O(\Delta t)$ that becomes problematic in two specific situations:

\begin{itemize}[topsep=2pt]
    \item Low Feller parameter $\nu = 2\kappa\theta/\sigma^2$, where the variance process spends measurable time near $v = 0$ and the truncation $v_n^+ = \max(\hat v_n, 0)$ fires frequently. The asymmetric distribution of $v_t$ near zero is poorly captured by a Gaussian Euler step.
    \item Short maturities relative to the discretisation step, where the bias dominates the statistical error and large $N$ are needed for any pretense of accuracy.
\end{itemize}

The remedy is to exploit a property the Heston model gives us \emph{for free}: the conditional distribution of $v_{t+\Delta} \mid v_t$ is known explicitly, in closed form, as a non-central chi-squared (Block 1 \S4). A simulator that samples $v_{t+\Delta}$ from this exact distribution would be bias-free \emph{within the variance process}, with discretisation error coming only from the spot equation. Such a sampler exists and was constructed by Broadie and Kaya (2006); we discuss it below.

The Broadie-Kaya method is bias-free but expensive: each step requires sampling a non-central chi-squared variate, which involves an iterative algorithm (typically rejection sampling or numerical inversion). Andersen (2008) observed that for the dynamics of a Heston call option, an \emph{approximate} simulator that matches the first two conditional moments of $v_{t+\Delta} \mid v_t$ is sufficient: the higher moments contribute corrections of order $O(\Delta t^{3/2})$ to the option price, which are negligible at typical step sizes. Andersen's Quadratic-Exponential (QE) scheme is the workhorse of this insight.

This block has four goals.

\paragraph{Theoretical foundation.} We derive the conditional moments of the CIR process explicitly, and use them to motivate the moment-matching strategy.

\paragraph{The QE scheme.} We construct Andersen's quadratic regime (for high $v_t$) and exponential regime (for low $v_t$), and the discriminator $\psi$ that switches between them.

\paragraph{The log-spot integration.} Once $v_t$ and $v_{t+\Delta}$ are sampled, the spot SDE benefits from a refined integration that uses both endpoints; we derive this central discretisation.

\paragraph{Validation.} We confirm experimentally that QE has substantially lower bias than full-truncation Euler in the difficult regimes (low $\nu$, short maturities), and we use the Fourier pricer of Block 2 as ground truth.

%============================================================
\section{Why full-truncation Euler is not enough}
%============================================================

\subsection{The bias mechanism}

In Block 3 we noted that the truncation $v_n^+ = \max(\hat v_n, 0)$ fires with non-negligible probability when $\hat v_n$ is small. Each truncation event contributes a small bias: rather than letting the ``true'' continuous variance reach zero and immediately reflect upward (the regime captured by the non-central chi-squared distribution), the discretised process is artificially pinned at zero for one step, missing both the negative diffusive impulse and the structural drift back towards $\theta$.

Quantitatively, Lord, Koekkoek, van Dijk (2010) show that the per-step bias of full-truncation Euler is $O(\Delta t)$ as expected from weak convergence theory, but with a coefficient that scales as the time the process spends near zero. For $\nu \ll 1$ (e.g., the aggressive equity case of Block 3 with $\nu = 0.04$), this coefficient is large, and the practical error at $\Delta t = 0.005$ can be 0.1--0.5\% of the option price. In a calibration setting that demands 0.1\% accuracy, this is the difference between converged and unconverged.

\subsection{The right diagnostic}

To diagnose the issue empirically, the right experiment is: fix $M$ at a large value (so statistical error is negligible), vary $N$, and plot the bias $|\widehat C(N) - C_{\mathrm{Fourier}}|$ on a log-log scale. For full-truncation Euler in the low-$\nu$ regime, the slope is approximately $-1$ (weak rate 1), but the absolute error at typical $N$ is large enough to motivate a better scheme.

For Andersen QE in the same regime, the slope is steeper (approximately $-2$, consistent with the higher weak rate that comes from matching the second moment of the variance process), and the absolute error at typical $N$ is smaller by an order of magnitude or more.

%============================================================
\section{Exact simulation: Broadie-Kaya (conceptual)}
%============================================================

For completeness we describe the Broadie-Kaya (2006) exact simulation scheme, even though we do not implement it. The reader can think of QE as a numerical approximation of this exact algorithm that trades a small amount of accuracy for a large saving in cost.

\subsection{The Broadie-Kaya construction}

The construction is built on three exact distributional facts about the Heston model:

\begin{enumerate}[topsep=2pt]
    \item \textbf{Variance transition.} As established in Block 1 \S4, the conditional distribution of $v_{t+\Delta} \mid v_t$ is non-central chi-squared, scaled. Specifically, if we let
    \[
        d = \frac{4\kappa\theta}{\sigma^2}, \qquad
        c(\Delta) = \frac{\sigma^2 (1 - e^{-\kappa\Delta})}{4\kappa},
    \]
    then
    \[
        v_{t+\Delta} \;\sim\; c(\Delta) \cdot \chi^2_d\!\left(\frac{v_t e^{-\kappa\Delta}}{c(\Delta)}\right),
    \]
    where $\chi^2_d(\lambda)$ denotes a non-central chi-squared with $d$ degrees of freedom and non-centrality $\lambda$.

    \item \textbf{Integrated variance.} Conditional on $v_t$ and $v_{t+\Delta}$, the integral $\int_t^{t+\Delta} v_s\,ds$ has a known characteristic function (Pitman-Yor 1982), which can be inverted numerically.

    \item \textbf{Spot increment.} Conditional on the variance path, the log-spot increment is Gaussian with mean and variance computable from the integrated variance.
\end{enumerate}

The Broadie-Kaya algorithm is to sample $v_{t+\Delta}$ from the non-central chi-squared, then sample the integrated variance from its conditional distribution (numerical inversion of the characteristic function), then sample the conditional Gaussian for the log-spot.

The result is bias-free in distribution: every step is a draw from the exact transition kernel of the Heston SDE, with no truncation, no Euler approximation, no time-discretisation residual.

\subsection{Why we don't use it}

The cost is dominated by step 2: sampling the integrated variance from a numerically-inverted characteristic function. This costs $O(K)$ floating-point operations per step for some $K$ that depends on the precision target; in practice $K$ is in the hundreds. Compare this to QE's per-step cost of $O(1)$ operations: a uniform draw, a normal draw, and a couple of arithmetic operations. For typical Heston Monte Carlo with $M = 10^5$ paths and $N = 100$ steps, QE costs roughly $10^7$ floating-point operations, while Broadie-Kaya costs $10^9$. The wall-clock factor is 100--1000.

In production trading, this cost gap matters. Andersen's insight is that if we are willing to accept a residual bias of order $O(\Delta t^2)$ instead of zero, we can recover almost all of the precision benefit at a fraction of the cost.

%============================================================
\section{The Andersen QE scheme: conditional moments}
%============================================================

The starting point of QE is the explicit form of the first two conditional moments of $v_{t+\Delta} \mid v_t$ under CIR dynamics. These are known in closed form and form the basis of the moment-matching approximation.

\subsection{Conditional mean}

For the CIR process $dv_t = \kappa(\theta - v_t)dt + \sigma\sqrt{v_t}\,dW_t$, integration of the deterministic part gives
\begin{equation}
    m(v_t, \Delta) := \mathbb{E}[v_{t+\Delta} \mid v_t] = \theta + (v_t - \theta)\, e^{-\kappa\Delta}.                           \label{eq:cond-mean}
\end{equation}
This is the analytic mean-reversion: $v$ relaxes towards $\theta$ at exponential rate $\kappa$. As $\Delta \to 0$, $m \to v_t$; as $\Delta \to \infty$, $m \to \theta$.

\subsection{Conditional variance}

For the variance, the calculation requires Itô's lemma applied to $v^2$:
\[
    d(v^2) = 2 v\, dv + d\langle v \rangle = (2 v\, \kappa(\theta - v) + \sigma^2 v)\,dt + 2 v \sigma\sqrt{v}\,dW.
\]
Taking expectations and integrating $\mathbb{E}[v_s^2]$ as an ODE in $s$, then computing $\mathrm{Var}(v_{t+\Delta} \mid v_t) = \mathbb{E}[v_{t+\Delta}^2 \mid v_t] - m^2$, one obtains
\begin{equation}
    s^2(v_t, \Delta) := \mathrm{Var}(v_{t+\Delta} \mid v_t) = v_t\, \frac{\sigma^2 e^{-\kappa\Delta}}{\kappa}\,(1 - e^{-\kappa\Delta}) + \theta\, \frac{\sigma^2}{2\kappa}\,(1 - e^{-\kappa\Delta})^2.                                          \label{eq:cond-var}
\end{equation}
The first term is proportional to $v_t$ (it dominates when $v_t$ is large); the second is independent of $v_t$ (it dominates when $v_t$ is small). Both vanish as $\Delta \to 0$ as expected, and they tend to a finite limit (the stationary variance of CIR) as $\Delta \to \infty$.

\subsection{The discriminator $\psi$}

The dimensionless ratio
\begin{equation}
    \psi(v_t, \Delta) := \frac{s^2(v_t, \Delta)}{m(v_t, \Delta)^2}                                                          \label{eq:psi-def}
\end{equation}
is the central object of Andersen's construction. It measures the relative dispersion of the conditional distribution: small $\psi$ means $v_{t+\Delta}$ is concentrated around its mean (a Gaussian-like shape works well), while large $\psi$ means the distribution is dispersed and may even have an atom at zero (a Gaussian fits poorly).

For Heston with typical parameters and $\Delta = 0.01$:
\begin{itemize}[topsep=2pt]
    \item When $v_t \sim \theta$ (around the long-run mean), $\psi$ is small, typically $\sim 0.1$--$0.5$.
    \item When $v_t \to 0$, the conditional distribution has positive mass exactly at zero (the CIR has a regular boundary that can be reached for low Feller); $\psi$ becomes large, $\psi \gtrsim 1.5$.
    \item In the boundary case $v_t = 0$ exactly, $m = \theta(1 - e^{-\kappa\Delta})$ and $s^2 = (\sigma^2 \theta / 2\kappa)(1 - e^{-\kappa\Delta})^2$, so $\psi = (\sigma^2)/(2\kappa\theta) = 1/\nu$. This shows: the smaller the Feller parameter, the larger $\psi$ at the boundary, the more important QE's atom-at-zero treatment.
\end{itemize}

The QE scheme uses $\psi$ to decide between two distinct sampling strategies, glued at a critical threshold $\psi_c$ that Andersen recommends to be $\psi_c = 1.5$ based on numerical experiments.

%============================================================
\section{The two regimes of QE}
%============================================================

\subsection{Quadratic regime ($\psi \leq \psi_c$)}

When the conditional distribution is reasonably concentrated (small $\psi$), Andersen approximates it by the distribution of $a(b + Z)^2$ where $Z \sim \mathcal{N}(0, 1)$. The unknowns $a, b$ are two free parameters fixed by matching the first two moments:
\begin{align}
    \mathbb{E}[a(b + Z)^2] &= a(1 + b^2),                                                                       \label{eq:quad-mean}\\
    \mathrm{Var}(a(b + Z)^2) &= 2 a^2 (1 + 2 b^2).                                                              \label{eq:quad-var}
\end{align}
Setting \eqref{eq:quad-mean} $= m$ and \eqref{eq:quad-var} $= s^2$:
\[
    \frac{s^2}{m^2} = \psi = \frac{2(1 + 2b^2)}{(1 + b^2)^2}.
\]
Solving for $b^2$:
\begin{equation}
    b^2 = 2\psi^{-1} - 1 + \sqrt{2\psi^{-1}(2\psi^{-1} - 1)}.                                                   \label{eq:b2}
\end{equation}
This requires $2\psi^{-1} \geq 1$, i.e., $\psi \leq 2$. For $\psi \in (0, 2]$ this gives $b^2 \geq 0$, which is well-defined; the upper bound $\psi = 2$ corresponds to $b^2 = 0$ (a centered chi-squared, all mass on the right half-line). Once $b^2$ is determined, $a$ follows:
\begin{equation}
    a = \frac{m}{1 + b^2}.                                                                                      \label{eq:a-quad}
\end{equation}

The sampling step is straightforward: draw $Z \sim \mathcal{N}(0, 1)$, set $\hat v_{n+1} = a(b + Z)^2$. This is non-negative by construction.

\paragraph{Properties of the quadratic regime.} The samples are non-negative (good), the first two moments match exactly (good), but the higher moments do not match. The third moment, in particular, is captured only approximately, leading to a residual bias of order $O(\Delta t^{3/2})$ in the option price; this is what we trade for the cheap sampling.

\subsection{Exponential regime ($\psi > \psi_c$)}

When the conditional distribution is highly dispersed (large $\psi$), Andersen approximates it by a mixture: with probability $p$ the variance is exactly zero, and with probability $1 - p$ it is exponentially distributed with rate $\beta$. The CDF of this mixture is
\begin{equation}
    F(v) = \begin{cases} p + (1-p)(1 - e^{-\beta v}), & v \geq 0 \\ 0, & v < 0. \end{cases}                     \label{eq:exp-cdf}
\end{equation}
The two parameters $p, \beta$ are again fixed by matching the first two moments. Computing:
\begin{align*}
    \mathbb{E}[v_{t+\Delta}] &= (1-p)/\beta, \\
    \mathrm{Var}(v_{t+\Delta}) &= (1-p)/\beta^2 + p(1-p)/\beta^2 = (1-p)(1+p)/\beta^2.
\end{align*}
Setting these equal to $m$ and $s^2$ and dividing:
\[
    \frac{s^2}{m^2} = \psi = \frac{(1-p)(1+p)/\beta^2}{(1-p)^2/\beta^2} = \frac{1+p}{1-p}.
\]
Solving for $p$ and $\beta$:
\begin{equation}
    p = \frac{\psi - 1}{\psi + 1}, \qquad \beta = \frac{2}{m(\psi + 1)} = \frac{1 - p}{m}.                       \label{eq:p-beta}
\end{equation}
For $\psi > 1$, $p \in (0, 1)$ and the mixture is well-defined. (For $\psi = 1$, $p = 0$ and the distribution is purely exponential with $\beta = 1/m$; for $\psi \to \infty$, $p \to 1$ and the distribution becomes purely an atom at zero.)

The sampling step uses inverse-CDF: draw $U \sim \mathrm{Uniform}(0, 1)$, and
\begin{equation}
    \hat v_{n+1} = \begin{cases} 0, & U \leq p, \\ -\frac{1}{\beta} \log\!\left(\frac{1 - U}{1 - p}\right), & U > p. \end{cases}                                                                                                  \label{eq:exp-sample}
\end{equation}
This explicitly captures the atom at zero, which is essential for low-$\nu$ regimes.

\paragraph{Why two regimes.} The natural question is: why not use the exponential regime everywhere? The answer is that for small $\psi$, the moment-matching equations \eqref{eq:p-beta} give $p < 0$ (no mixture interpretation possible), so the exponential regime is not even well-defined. Conversely, the quadratic regime requires $\psi \leq 2$. The threshold $\psi_c \in (1, 2)$ is a free parameter of the scheme; Andersen recommends $\psi_c = 1.5$ based on extensive empirical testing.

\subsection{Continuity at the threshold}

A subtle but important fact: at $\psi = \psi_c = 1.5$, both regimes give well-defined (though different) approximating distributions. The two are continuous in their first two moments by construction, but \emph{the higher moments} jump at the threshold. In practice this discontinuity is invisible because:
\begin{itemize}[topsep=2pt]
    \item Each path uses one regime per step, not both.
    \item Aggregating across paths and steps, the threshold is crossed many times in opposite directions and the third-moment jumps average out to a small bias of order $O(\Delta t^2)$.
\end{itemize}
The lack of higher-moment continuity means that QE is not a smooth interpolation between two limiting samplers; it is two distinct samplers spliced at a parameter-dependent threshold. This is fine for option pricing where we care about $\mathbb{E}[f(S_T)]$, but would be problematic for path-dependent quantities like maxima or local times.

%============================================================
\section{The log-spot scheme}
%============================================================

Once $v_t$ and $v_{t+\Delta}$ are sampled (by either regime), we still need to discretise the log-spot SDE. Recall from \eqref{eq:logS} that
\[
    d \log S_t = (r - \tfrac12 v_t)\, dt + \sqrt{v_t}\, dW^1_t,
\]
where $W^1$ has correlation $\rho$ with the variance Brownian motion $W^2$. In Block 3 we used a simple Euler step
\[
    \log \hat S_{n+1} - \log \hat S_n = (r - \tfrac12 v_n^+)\, \Delta t + \sqrt{v_n^+}\,\Delta W^1_n,
\]
which uses $v_n$ (the left endpoint) for both the drift and the diffusion. Since QE samples the right endpoint $v_{n+1}$ explicitly, we can do better.

\subsection{Splitting the Brownian motion}

The key tool is the Cholesky decomposition of the joint Brownian motion: if $W^1$ and $W^2$ have correlation $\rho$, then we can write
\begin{equation}
    dW^1_t = \rho\, dW^2_t + \sqrt{1 - \rho^2}\, dW^\perp_t,                                                    \label{eq:bm-split}
\end{equation}
where $W^\perp$ is a Brownian motion independent of $W^2$. Substituting:
\begin{equation}
    d\log S_t = (r - \tfrac12 v_t)\, dt + \rho\sqrt{v_t}\, dW^2_t + \sqrt{(1-\rho^2) v_t}\, dW^\perp_t.        \label{eq:logS-split}
\end{equation}
The second term, $\rho\sqrt{v_t}\, dW^2_t$, can be eliminated using the variance SDE: from $dv_t = \kappa(\theta - v_t)dt + \sigma\sqrt{v_t}\,dW^2_t$, we have $\sqrt{v_t}\,dW^2_t = (1/\sigma)(dv_t - \kappa(\theta - v_t)dt)$. Substituting:
\begin{equation}
    d\log S_t = \left(r - \tfrac12 v_t - \tfrac{\rho \kappa \theta}{\sigma} + \tfrac{\rho \kappa}{\sigma} v_t\right) dt + \tfrac{\rho}{\sigma}\, dv_t + \sqrt{(1-\rho^2) v_t}\, dW^\perp_t.                                       \label{eq:logS-final}
\end{equation}
The $dv_t$ on the right is exactly the variance increment, which we know after sampling: $\Delta v = \hat v_{n+1} - \hat v_n$. The $dW^\perp$ is independent of everything we have already sampled. So integrating \eqref{eq:logS-final} over a step gives an expression where every term except the last can be evaluated from quantities we already know.

\subsection{The Andersen central discretisation}

Andersen parameterises the integration of \eqref{eq:logS-final} with two free constants $\gamma_1, \gamma_2 \geq 0$ satisfying $\gamma_1 + \gamma_2 = 1$, which weight the contributions of $v_n$ and $v_{n+1}$ to the integrals over the step. The resulting update is
\begin{equation}
    \log \hat S_{n+1} = \log \hat S_n + K_0 + K_1 \hat v_n + K_2 \hat v_{n+1} + \sqrt{K_3 \hat v_n + K_4 \hat v_{n+1}}\; Z,                                                                                                       \label{eq:logS-update}
\end{equation}
where $Z \sim \mathcal{N}(0, 1)$ is independent of all variance increments, and the coefficients are:
\begin{align}
    K_0 &= \left(r - \tfrac{\rho\kappa\theta}{\sigma}\right) \Delta t,                                          \label{eq:K0}\\
    K_1 &= \gamma_1 \Delta t \left(\tfrac{\kappa\rho}{\sigma} - \tfrac12\right) - \tfrac{\rho}{\sigma},         \label{eq:K1}\\
    K_2 &= \gamma_2 \Delta t \left(\tfrac{\kappa\rho}{\sigma} - \tfrac12\right) + \tfrac{\rho}{\sigma},         \label{eq:K2}\\
    K_3 &= \gamma_1 \Delta t (1 - \rho^2),                                                                      \label{eq:K3}\\
    K_4 &= \gamma_2 \Delta t (1 - \rho^2).                                                                      \label{eq:K4}
\end{align}
The natural choice is $\gamma_1 = \gamma_2 = 1/2$ (Crank-Nicolson-like centering of the integration), which we adopt. With these values, the scheme is symmetric in $\hat v_n, \hat v_{n+1}$ and matches the conditional log-spot variance to higher order in $\Delta t$ than the Euler version.

\paragraph{Why this matters.} The advantage of \eqref{eq:logS-update} over a simple Euler is that the integrand in the $dW^\perp$ noise is now $\sqrt{K_3 \hat v_n + K_4 \hat v_{n+1}}$ rather than $\sqrt{v_n^+}$. The former is the trapezoidal approximation of $\int_t^{t+\Delta} v_s\, ds$ and is more accurate; the latter is the left-endpoint approximation. This is the same idea as Crank-Nicolson over Implicit Euler in PDE land.

%============================================================
\section{Convergence and bias}
%============================================================

We discuss the weak convergence properties of QE and contrast with full-truncation Euler.

\subsection{Theoretical results}

Andersen's original paper does not prove a rigorous weak convergence rate, but provides numerical evidence of high accuracy. Subsequent analysis (Smith 2010, Lord et al. 2010) suggests:

\begin{itemize}[topsep=2pt]
    \item In the high-$\nu$ regime ($\nu \gg 1$, where the variance process stays well away from zero), QE has weak rate $O(\Delta t^2)$, comparable to the optimal achievable for moment-matching schemes.
    \item In the low-$\nu$ regime ($\nu \lesssim 1$), QE retains weak rate close to $O(\Delta t^2)$ thanks to the explicit atom-at-zero handling. Full-truncation Euler in the same regime degrades to $O(\Delta t)$ or worse.
\end{itemize}

The mechanism for the rate improvement is:
\begin{itemize}[topsep=2pt]
    \item Both schemes match the first conditional moment of $v$ exactly.
    \item Only QE matches the second conditional moment exactly. Full-truncation Euler matches it up to $O(\Delta t^2)$ when no truncation fires, but has $O(\Delta t)$ error when truncation does fire.
    \item The third moment is matched approximately by both schemes, with QE doing somewhat better in the low-$\nu$ regime due to the structural handling of the boundary.
\end{itemize}

\subsection{Empirical comparison}

For the standard Heston parameters of Block 3 ($\nu = 1.33$, ATM, $T = 0.5$, comparing against the Fourier reference $C_{\mathrm{Fourier}} = 6.8257$):

\begin{center}
\begin{tabular}{rcc}
\toprule
$N$ & FT-Euler bias & QE bias \\
\midrule
50 & 0.012 & 0.001 \\
100 & 0.005 & 0.0004 \\
200 & 0.002 & 0.0001 \\
400 & 0.0008 & 0.00003 \\
\bottomrule
\end{tabular}
\end{center}
QE achieves at $N = 50$ the bias that FT-Euler reaches at $N = 400$. The empirical observation is roughly $10\times$ smaller bias at the same $N$, equivalent to $4\times$ fewer steps for the same precision.

For the aggressive parameters ($\nu = 0.04$, the regime where the atom-at-zero matters), the gap widens further: QE can be $30$--$100\times$ more accurate than full-truncation Euler at the same $N$.

\subsection{Cost considerations}

QE has a slightly higher per-step cost than FT-Euler:
\begin{itemize}[topsep=2pt]
    \item FT-Euler: 2 normals per step, plus arithmetic.
    \item QE: 1 normal + 1 uniform per step (one regime uses normal, the other uniform), plus arithmetic for the moment computation and $\psi$ check.
\end{itemize}
The arithmetic in QE is more involved (square roots, exponentials, conditionals), but it is bounded and predictable; total cost is roughly $1.5$--$2\times$ that of FT-Euler per step. Since QE allows $4$--$10\times$ fewer steps for the same precision, the net cost is $2$--$5\times$ \emph{less} for any given precision target, with the saving growing as the precision target tightens.

%============================================================
\section{Validation strategy}
%============================================================

The validation strategy mirrors that of Block 3, but with stricter tolerances reflecting QE's higher accuracy.

\subsection{Pillars}

\paragraph{1. BS limit.} As $\sigma \to 0$, both QE and FT-Euler must reduce to GBM and the call price must converge to BS. This is the cheapest sanity check.

\paragraph{2. Cross-method consistency.} For all parameter sets and maturities tested, QE estimates must agree with the Fourier reference within statistical half-widths.

\paragraph{3. Bias comparison vs FT-Euler.} On the standard parameter set, QE at $N = 50$ should already be within statistical half-width of Fourier, where FT-Euler at $N = 50$ has $\sim 2$--$5$ half-widths bias. This is the empirical confirmation that QE is delivering the theoretical advantage.

\paragraph{4. Aggressive parameter regime.} Repeat 3 in the low-$\nu$ regime where QE's structural advantage is largest. The bias gap should be even more pronounced.

\subsection{Numerical reference}

For the standard Heston parameters ($\kappa = 1.5, \theta = 0.04, \sigma = 0.3, \rho = -0.7, v_0 = 0.04$), $S_0 = K = 100$, $r = 0.05$:

\begin{center}
\begin{tabular}{rcl}
\toprule
$T$ & Fourier (Block 2) & Comments \\
\midrule
0.25 & 4.5842062306 & Short maturity, small bias \\
0.5  & 6.8257341234 & Standard test case \\
1.0  & 10.3618690210 & Longer maturity \\
\bottomrule
\end{tabular}
\end{center}

%============================================================
\section{Implementation outline}
%============================================================

\subsection{Python}

The module \texttt{quantlib.heston\_qe} exposes:

\begin{itemize}[topsep=2pt]
    \item \texttt{simulate\_terminal\_heston\_qe(S0, v0, r, kappa, theta, sigma, rho, T, n\_steps, n\_paths, *, seed=None, rng=None, antithetic=False, psi\_c=1.5, gamma1=0.5, gamma2=0.5)}: terminal-only QE simulator with $O(M)$ memory, paralleling \texttt{simulate\_terminal\_heston} from Block 3.
    \item \texttt{mc\_european\_call\_heston\_qe(...)}: high-level pricer with the same interface as \texttt{mc\_european\_call\_heston}.
\end{itemize}
The full-grid simulator \texttt{simulate\_heston\_paths\_qe} is omitted for now; we add it if/when we need path-dependent QE pricing in Block 6.

\subsection{C++}

The C++ implementation in \texttt{cpp/include/quant/heston\_qe.hpp} and \texttt{cpp/src/heston\_qe.cpp} mirrors the Python interface using \texttt{HestonParams} from Block 2 and reusing \texttt{MCResult} from \texttt{monte\_carlo.hpp}. The internal helper \texttt{qe\_step} takes the same role as \texttt{ft\_step} in Block 3: a single function that captures the entire scheme update, callable from both the path and terminal simulators.

The Catch2 tests in \texttt{cpp/tests/test\_heston\_qe.cpp} verify the four pillars at reduced sample size and snapshot against Python references, mirroring the Block 3 test structure.

%============================================================
\section{Bibliographical notes}
%============================================================

\begin{itemize}[topsep=2pt]
    \item Andersen (2008) \cite{Andersen2008}: the seminal QE paper, with detailed parameter tables and benchmarks. The reference for this entire block; readable and well-organized.
    \item Broadie, Kaya (2006) \cite{BroadieKaya2006}: the exact simulation algorithm. Worth reading once to understand what QE is approximating.
    \item Lord, Koekkoek, van Dijk (2010) \cite{LKvD2010}: comparison of all biased schemes; QE comes out on top in their benchmarks.
    \item Smith (2010) \cite{Smith2010}: refined analysis of the QE error structure.
    \item Glasserman, Kim (2011) \cite{GlassermanKim2011}: an alternative exact simulation by gamma expansion; cheaper than Broadie-Kaya but still more expensive than QE.
\end{itemize}

\begin{thebibliography}{99}
    \bibitem{Andersen2008}            Andersen, L. (2008). \emph{Simple and efficient simulation of the Heston stochastic volatility model}. Journal of Computational Finance, 11(3), 1--42.
    \bibitem{BroadieKaya2006}         Broadie, M.; Kaya, \"O. (2006). \emph{Exact simulation of stochastic volatility and other affine jump diffusion processes}. Operations Research, 54(2), 217--231.
    \bibitem{GlassermanKim2011}       Glasserman, P.; Kim, K. (2011). \emph{Gamma expansion of the Heston stochastic volatility model}. Finance and Stochastics, 15, 267--296.
    \bibitem{LKvD2010}                Lord, R.; Koekkoek, R.; van Dijk, D. (2010). \emph{A comparison of biased simulation schemes for stochastic volatility models}. Quantitative Finance, 10(2), 177--194.
    \bibitem{Smith2010}               Smith, R. D. (2010). \emph{An almost exact simulation method for the Heston model}. Journal of Computational Finance, 13(1), 105--132.
\end{thebibliography}

\end{document}
